% Modality Performance Subsection
\subsection{Modality Performance}

Biosignal modality selection differs substantially between detection and forecasting approaches. Detection studies prioritize movement sensors, while forecasting studies emphasize autonomic nervous system indicators.

\subsubsection{Modality Usage}

Accelerometer data appears in nine of 13 studies (69\%). Six detection studies and three forecasting studies use accelerometry. Electrodermal activity appears in five studies (38\%). ECG or HRV data appears in five studies. PPG appears in four studies. Gyroscope data appears in three studies, all in detection applications.

Multi-modal approaches dominate the evidence base. Nine of 13 studies (69\%) use multiple sensor types. Four studies use single-modality approaches: \textcite{spahrDeepLearningbasedDetection2025} (ACC only), \textcite{reintjesECGBasedDetectionEpileptic2025} (ECG only), \textcite{borujenyDetectionEpilepticSeizure2013} (ACC only), and \textcite{odeDevelopmentEpilepticSeizure2023} (ECG RRI only).

\subsubsection{Sensor Location}

Wrist-worn devices appear in eight studies (62\%). Arm or armband placement appears in three studies. Chest placement appears in one study. Thigh placement appears in one study. \textcite{borujenyDetectionEpilepticSeizure2013} used multiple sensor locations (right arm, left arm, left thigh).

The preference for wrist placement reflects patient acceptance and the availability of consumer wearable devices. The Empatica E4 wristband appears in three studies. The Embrace watch appears in one study. Fitbit smartwatches appear in one study.

\subsubsection{Detection Performance by Modality}

Single-modality accelerometer approaches achieve strong performance. \textcite{spahrDeepLearningbasedDetection2025} reported 96\% sensitivity with 0.005/h FPR using only ACC data from the Empatica E4. This represents the lowest FPR among all detection studies. \textcite{borujenyDetectionEpilepticSeizure2013} reported 85\% sensitivity using 2D ACC, though with only three patients.

Multi-modal movement detection shows excellent sensitivity in controlled settings. \textcite{fineWearableDeviceSeizure2025} used a 6-axis band combining ACC and gyroscope to achieve 100\% sensitivity with 0.023/h FPR in an independent test set. However, the test set contained only 10 seizures from 3 patients.

ECG-based detection shows high FPR variability. \textcite{reintjesECGBasedDetectionEpileptic2025} compared three anomaly detection methods on single-lead ECG. Sensitivity ranged from 38.0\% to 98.2\%, while FPR ranged from 1.91/h to 39.75/h depending on the method. \textcite{odeDevelopmentEpilepticSeizure2023} achieved 74\% sensitivity with 0.85/h FPR using ECG RRI data.

Multi-modal approaches combining movement and autonomic signals show mixed results. \textcite{dongTwoLayerEnsembleMethod2022} combined ACC and PPG in a two-stage CNN-LSTM architecture, achieving 71.6\% sensitivity with 0.165/h FPR. \textcite{wangDevelopmentWearableSeizure2025} combined ACC, gyroscope, SEMG, and EDA to achieve 56.4\% to 95.3\% sensitivity depending on the feature set.

\textcite{wangDevelopmentWearableSeizure2025} directly compared modality performance. Attitude angle features (PITCH and ROLL) outperformed raw accelerometer and gyroscope data for seizure detection. This finding suggests that feature engineering can enhance the value of standard movement sensors.

\subsubsection{Forecasting Performance by Modality}

Forecasting studies consistently use autonomic signals. \textcite{nasseriAmbulatorySeizureForecasting2021} combined ACC, PPG, EDA, temperature, and HR to achieve AUC 0.75 (SD 0.15). \textcite{meiselMachineLearningWristband2020} used EDA, PPG, temperature, and ACC from the Empatica E4 to achieve 51.2\% sensitivity with IoC 14.1\% in LOSO validation.

\textcite{vielufSeizureMonitoringCombined2025} combined EDA, ACC, and temperature with diary data to achieve 82\% sensitivity and 67\% specificity. This study demonstrated that 24-hour harmonic patterns in EDA and accelerometry contain predictive information.

\textcite{stirlingForecastingSeizureLikelihood2021} focused on heart rate cycles combined with sleep and step data from Fitbit devices. The approach achieved AUC 0.74 with 100\% patient success in hourly forecasting mode.

\textbf{Key finding.} Multi-modal approaches are common but not universally superior. Single-modality ACC (\textcite{spahrDeepLearningbasedDetection2025}) achieves the best FPR in the detection literature. Forecasting consistently requires multi-modal autonomic data, but AUC plateaus around 0.75 regardless of modality combination.
